%% toc: P3 40 40.5
\chapter{Closure Constructions and the Dynamic View of Sparse Random Graphs}
\label{sec:p03-w03-p03-068-closure}

In the sparse random graph \(G(n,n^{-\alpha})\) with \(\alpha\in(0,1)\)
irrational, a bounded set of vertices can have special properties that a
random set of the same size almost surely lacks: three vertices may, for
instance, have a common neighbour although three random vertices almost surely
do not (\cref{def:p03-w03-p03-065-rigid-special}).  Spencer's closure
construction collects, around a finite vertex set \(X\), exactly the vertices
that witness such special properties: the vertices reachable from \(X\) by
repeatedly adding rigid extensions of bounded size.  This chapter gives the
two definitions of the \(t\)-closure \(\mathrm{cl}_t(X)\), proves that they
are well defined and agree, shows that the isomorphism type of a closure is a
first-order property of the base tuple, and proves the two probabilistic
facts about closures that the analysis of \(G(n,n^{-\alpha})\) needs: almost
surely the closure of the empty set is empty, and almost surely every
\(r\)-tuple has a closure of bounded size.  The second half of the chapter
fixes a first-order sentence and lets \(\alpha\) vary.  The limiting
probability is then a \(\{0,1\}\)-valued function of the irrational
exponent; it is locally constant, its discontinuities can lie only at
rationals, and every finite set of rationals in \((0,1)\) is the set of
discontinuities of some sentence.

\paragraph{Inputs and conclusions.}
The inputs are a graph \(G\) (for the deterministic part an arbitrary
graph, for the probabilistic part the random graph \(G(n,n^{-\alpha})\) of
\cref{def:p03-w03-p03-065-gnp}), a set or tuple of its vertices, an
irrational \(\alpha\in(0,1)\) and a bound \(t\geq0\) on the size of the rigid
extensions used.  The conclusions are: the dichotomy that an extension over
a finite root set is safe or has a rigid part
(\cref{thm:p03-w03-p03-068-safe-or-rigid}); the closure \(\mathrm{cl}_t(X)\) and
its order-independence (\cref{thm:p03-w03-p03-068-well-defined}); a
first-order formula expressing each closure type
(\cref{thm:p03-w03-p03-068-type-formula}); almost sure emptiness of
\(\mathrm{cl}_t(\emptyset)\) (\cref{thm:p03-w03-p03-068-nonexistence}); an
almost sure bound \(r+K\) on the size of the \(t\)-closure of every
\(r\)-tuple (\cref{thm:p03-w03-p03-068-finite-closure}); and, for a sentence
\(A\), the local constancy of \(\alpha\mapsto\lim_n\Pr[G(n,n^{-\alpha})\entails
A]\) (\cref{thm:p03-w03-p03-068-continuity}) together with the description
of its finite spectra (\cref{thm:p03-w03-p03-068-finite-spectra}).

\paragraph{Selected treatment.}
The presentation follows Joel Spencer's \emph{Logic and Random Structures}
in the author's version hosted on his New York University page
\cite{Spencer2007LogicRandomStructuresAuthor}, not the printed chapter in the 2007 Springer
volume \emph{Finite Model Theory and Its Applications}, which was not
consulted.  All locators refer to the author's version: closure is its Section~4
\cite[pp.~8--9]{Spencer2007LogicRandomStructuresAuthor} and the dynamic view its Section~8
\cite[pp.~15--16]{Spencer2007LogicRandomStructuresAuthor}.  A shorter
account of the same closure material, with its own proofs, is in
\cref{sec:p03-w03-p03-064-closure}.  Rooted graphs, their parameters
\(r,v,e\) and the notions dense, sparse, rigid and safe are those of
\cref{def:p03-w03-p03-065-rooted-graph,def:p03-w03-p03-065-dense-sparse},
taken from the source's Section~3.  The source states that the
equivalence and well-definedness of its two closure definitions rest on
elementary combinatorial lemmas which it does not give
\cite[p.~8]{Spencer2007LogicRandomStructuresAuthor}; those lemmas are stated and proved
here (\cref{thm:p03-w03-p03-068-rigid-local,thm:p03-w03-p03-068-move-base}).
The proofs of the Nonexistence and Finite Closure Lemmas expand the
source's arguments \cite[pp.~8--9]{Spencer2007LogicRandomStructuresAuthor}.  The
continuity argument of the dynamic view is the source's
\cite[p.~15]{Spencer2007LogicRandomStructuresAuthor}, written out from the explicit axioms
of the almost sure theory.  The proofs that the almost sure theories are
pairwise distinct and that every finite set of rationals is a spectrum are
written for this chapter: the source states both facts, for the first
refers to a graph whose copies have a threshold at a given rational, and for
the second says only that some technical work suffices
\cite[pp.~15--16]{Spencer2007LogicRandomStructuresAuthor}.

\section{Rigid extensions over a vertex set}
\label{sec:p03-w03-p03-068-rigid}

\begin{hbconvention}{def:p03-w03-p03-068-standing}{Standing conventions}
A \emph{graph} is a set of vertices with an irreflexive symmetric adjacency
relation \(\sim\); it may be infinite unless it is said to be finite.  For a
set \(W\) of vertices of \(G\), \(G[W]\) is the induced subgraph on \(W\).
Throughout, \(\alpha\in(0,1)\) is a fixed irrational number and \(t\geq0\)
is an integer, as in Spencer's Section~4, which fixes \(\alpha\) and
\(t\geq1\) and adds the case \(t=0\) later
\cite[pp.~8--9]{Spencer2007LogicRandomStructuresAuthor}.  For a set \(S\) of vertices and a
finite nonempty set \(N\) of vertices disjoint from \(S\), let
\(e_G(N;S)\) be the number of edges of \(G\) with at least one end in \(N\)
and both ends in \(S\cup N\).  This is a cardinal; if it is infinite, the
expression \(|N|-e_G(N;S)\,\alpha\) is read as \(-\infty\).
\end{hbconvention}

In the language of rooted graphs, if \(S\) is finite then \(e_G(N;S)\) is
the parameter \(e\) of the rooted graph \((S,G[S\cup N])\), whose nonroots
are the \(v=|N|\) vertices of \(N\): edges between two roots are exactly the
ones not counted.

\begin{hbdefinition}{def:p03-w03-p03-068-rigid-over}{Extensions and rigid extensions over a vertex set}
Let \(G\) be a graph, let \(x_1,\ldots,x_r\) be distinct vertices, and let
\((R,H)\) be a rooted graph with roots \(x_1,\ldots,x_r\) and nonroots
\(y_1,\ldots,y_v\) in the labelling of
\cref{def:p03-w03-p03-065-rooted-graph}.  Distinct vertices
\(b_1,\ldots,b_v\) of \(G\), distinct from the \(x_i\), \emph{form an
\((R,H)\) extension over} \(x_1,\ldots,x_r\) if \(b_i\sim b_j\) whenever
\(y_iy_j\in E(H)\) and \(x_i\sim b_j\) whenever \(x_iy_j\in E(H)\); edges
between roots are not examined and further edges are allowed.  They
\emph{form a rigid extension over} \(x_1,\ldots,x_r\) if they form an
\((R,H)\) extension for some rigid \((R,H)\)
\cite[p.~8]{Spencer2007LogicRandomStructuresAuthor}.

For an arbitrary set \(X\) of vertices and a finite nonempty set \(Y\) of
vertices disjoint from \(X\), say that \(Y\) is \emph{rigid over} \(X\) if
\[
  |Y\setminus S|-e_G(Y\setminus S;S)\,\alpha<0
  \qquad\text{for every }S\text{ with }X\subseteq S\subsetneq X\cup Y.
\]
A set \(Z\) of vertices is \emph{\(t\)-closed} if no set \(Y\) with
\(1\leq|Y|\leq t\) and \(Y\cap Z=\emptyset\) is rigid over \(Z\).
\end{hbdefinition}

Taking \(S=X\) shows that a rigid \(Y\) satisfies
\(|Y|-e_G(Y;X)\alpha<0\): rigidity includes density, as for rooted graphs.
For \(t=0\) every set is \(t\)-closed, because there is no \(Y\) with
\(1\leq|Y|\leq0\).

\begin{hblemma}{thm:p03-w03-p03-068-rigid-local}{proved=annex:p03-w03-p03-068-rigid-local}{Rigidity is a property of the induced graph}
Let \(X=\{x_1,\ldots,x_r\}\) be a finite set of \(r\) distinct vertices of a
graph \(G\) and \(Y=\{b_1,\ldots,b_v\}\) a set of \(v\geq1\) vertices
disjoint from \(X\).  The following are equivalent.
\begin{enumerate}
\item \(b_1,\ldots,b_v\) form a rigid extension over \(x_1,\ldots,x_r\).
\item The rooted graph \((X,G[X\cup Y])\) is rigid.
\item \(Y\) is rigid over \(X\).
\end{enumerate}
Moreover, for any set \(W\supseteq X\cup Y\) of vertices, \(Y\) is rigid over
\(X\) in \(G\) if and only if it is rigid over \(X\) in \(G[W]\); and the same
holds for an infinite \(X\) whenever \(X\cup Y\subseteq W\).
\end{hblemma}

\begin{hblemma}{thm:p03-w03-p03-068-move-base}{proved=annex:p03-w03-p03-068-move-base}{Moving to a larger base}
Let \(X\subseteq B\) be sets of vertices of a graph \(G\), let \(Y\) be rigid
over \(X\), and suppose \(Y'=Y\setminus B\) is nonempty.  Then \(Y'\) is
rigid over \(B\), and \(|Y'|\leq|Y|\).
\end{hblemma}

The lemma says that a rigid extension never becomes less rigid when more
vertices are nailed down as roots and some of its own vertices are absorbed
into the base.  It is the combinatorial fact behind every property of
closures below.

Rigid extensions sit at one end of the dense--sparse spectrum; safe
extensions (\cref{def:p03-w03-p03-065-dense-sparse}) sit at the other.
Over a finite root set, a finite extension is either safe or has a nonempty
part that is rigid over the roots, and never both.  Say that \(b_1,\ldots,b_v\) \emph{form a safe
extension over} \(x_1,\ldots,x_r\) if the rooted graph
\((X,G[X\cup Y])\) is safe, with \(X\) and \(Y\) as in
\cref{thm:p03-w03-p03-068-rigid-local}.  Spencer's completeness argument
uses, with its proof omitted, the combinatorial fact that a nonsafe
extension contains a rigid subextension
\cite[p.~12]{Spencer2007LogicRandomStructuresAuthor}, and he speaks of
safe extensions when he describes closures in countable models
\cite[p.~15]{Spencer2007LogicRandomStructuresAuthor}.  The next lemma
states and proves that fact.

\begin{hblemma}{thm:p03-w03-p03-068-safe-or-rigid}{proved=annex:p03-w03-p03-068-safe-or-rigid}{Safe, or containing a rigid part}
Let \(X\) be a finite set of vertices of a graph \(G\), \(\alpha\in(0,1)\)
irrational, and \(Y\) a finite nonempty set of vertices disjoint from \(X\).
Then exactly one of the following holds.
\begin{enumerate}
\item The rooted graph \((X,G[X\cup Y])\) is safe.
\item Some nonempty \(U\subseteq Y\) is rigid over \(X\).
\end{enumerate}
Consequently a finite set \(Z\) is \(t\)-closed if and only if
\((Z,G[Z\cup Y])\) is safe for every \(Y\) with \(1\leq|Y|\leq t\) and
\(Y\cap Z=\emptyset\).
\end{hblemma}

For \(\alpha=\pi/7\) and \(X=\{x_1,x_2,x_3,x_4\}\), a vertex \(y_1\)
adjacent to all four together with a vertex \(y_2\) adjacent only to \(y_1\)
is neither safe nor rigid over \(X\) (Spencer's example of a dense rooted
graph that is not rigid \cite[p.~7]{Spencer2007LogicRandomStructuresAuthor});
the lemma finds the rigid part \(U=\{y_1\}\), with \(1-4\alpha<0\).

\section{The \texorpdfstring{\(t\)}{t}-closure}
\label{sec:p03-w03-p03-068-closure-def}

Spencer gives two definitions \cite[p.~8]{Spencer2007LogicRandomStructuresAuthor}.

\begin{hbdefinition}{def:p03-w03-p03-068-closure}{The \(t\)-closure: algorithmic and minimal-set definitions}
Let \(G\) be a graph and \(X\) a set of vertices.
\begin{enumerate}
\item \emph{Algorithmic.}  A \emph{rigid \(t\)-chain from \(X\)} is a
finite or infinite sequence \(X=X_0\subsetneq X_1\subsetneq\cdots\) with
\(X_{i+1}=X_i\cup Y_i\), where \(1\leq|Y_i|\leq t\), \(Y_i\cap
X_i=\emptyset\) and \(Y_i\) is rigid over \(X_i\).  A finite chain
\(X_0\subsetneq\cdots\subsetneq X_l\) is \emph{terminal} if \(X_l\) is
\(t\)-closed.  Spencer's procedure begins with \(X\), adds the vertices of
any rigid extension of at most \(t\) vertices over the current set, and
iterates until none is left; the final set is a terminal chain's last set.
\item \emph{Minimal set.}  \(\mathrm{cl}_t(X)\) is the least \(t\)-closed set
of vertices containing \(X\).
\end{enumerate}
For a tuple of vertices,
\(\mathrm{cl}_t(x_1,\ldots,x_r)=\mathrm{cl}_t(\{x_1,\ldots,x_r\})\).
\end{hbdefinition}

That the least set exists, that the procedure ends there whatever choices
it makes, and when it ends at all, is the content of the next theorem.
For \(t=0\) both definitions give \(\mathrm{cl}_0(X)=X\), which is
Spencer's convention for the \(0\)-closure \cite[p.~9]{Spencer2007LogicRandomStructuresAuthor}.

\begin{hbtheorem}{thm:p03-w03-p03-068-well-defined}{proved=annex:p03-w03-p03-068-well-defined}{The two definitions are well defined and agree}
Let \(G\) be a graph, \(X\) a set of vertices and \(t\geq0\).
\begin{enumerate}
\item The intersection of any nonempty family of \(t\)-closed sets is
\(t\)-closed, and the set of all vertices is \(t\)-closed.  Hence
\(\mathrm{cl}_t(X)\), the intersection of all \(t\)-closed sets containing
\(X\), is the least \(t\)-closed set containing \(X\).
\item Every set of every rigid \(t\)-chain from \(X\) is contained in
\(\mathrm{cl}_t(X)\).
\item The last set of every terminal chain from \(X\) equals
\(\mathrm{cl}_t(X)\).  In particular it does not depend on the order in
which rigid extensions are added or on which ones are chosen.
\item If \(\mathrm{cl}_t(X)\setminus X\) is finite, in particular if \(G\) is
finite, then every rigid \(t\)-chain from \(X\) has at most
\(|\mathrm{cl}_t(X)\setminus X|\) steps and extends to a terminal chain;
so the procedure always stops, and stops at \(\mathrm{cl}_t(X)\).
\item \(X\subseteq\mathrm{cl}_t(X)\);
\(X\subseteq X'\) implies \(\mathrm{cl}_t(X)\subseteq\mathrm{cl}_t(X')\);
\(s\leq t\) implies \(\mathrm{cl}_s(X)\subseteq\mathrm{cl}_t(X)\); and
\(\mathrm{cl}_t(\mathrm{cl}_t(X))=\mathrm{cl}_t(X)\).
\item If \(W\supseteq\mathrm{cl}_t(X)\), then the \(t\)-closure of \(X\)
computed in \(G[W]\) equals \(\mathrm{cl}_t(X)\).  In particular, the
\(t\)-closure of \(X\) computed in \(G[\mathrm{cl}_t(X)]\) is all of
\(\mathrm{cl}_t(X)\).
\end{enumerate}
\end{hbtheorem}

\begin{hbexample}{def:p03-w03-p03-068-example}{A closure at \(\alpha=\pi/7\)}
Spencer's example \cite[p.~8]{Spencer2007LogicRandomStructuresAuthor}: \(\mathrm{cl}_4(x_1,x_2)\)
might consist of \(x_1,x_2\); \(y_1,y_2\) adjacent to each other and to both
\(x_1,x_2\); \(y_3,y_4,y_5,y_6\) forming a \(K_5\) with \(y_2\); and \(y_7\), a
common neighbour of \(x_2,y_1,y_5\).  Take \(\alpha=\pi/7=0.4487\ldots\) and
let \(G\) be the graph on these nine vertices with exactly these \(8+10=18\)
edges (the eight listed edges at \(x_1,x_2,y_1,y_7\) and the ten edges of the
\(K_5\)).
\begin{itemize}
\item \(\{y_1,y_2\}\) is rigid over \(\{x_1,x_2\}\): with \(S=\{x_1,x_2\}\)
there are \(v=2\) new vertices and \(e=5\) counted edges, \(2-5\alpha<0\); with
one of \(y_1,y_2\) nailed down, \(v=1\), \(e=3\), \(1-3\alpha<0\).
\item \(\{y_3,y_4,y_5,y_6\}\) is rigid over \(\{x_1,x_2,y_1,y_2\}\): if \(k\)
of the four remain unnailed they carry \(\binom k2+k(5-k)=k(9-k)/2\) counted
edges, and \(k-\alpha k(9-k)/2<0\) for \(k=1,2,3,4\), the values being
\(1-4\alpha\), \(2-7\alpha\), \(3-9\alpha\), \(4-10\alpha\).
\item \(\{y_7\}\) is rigid over the remaining seven vertices: \(1-3\alpha<0\).
\end{itemize}
These three steps form a terminal chain, so \(\mathrm{cl}_4(x_1,x_2)=V(G)\).
By contrast \(\mathrm{cl}_3(x_1,x_2)=\mathrm{cl}_2(x_1,x_2)=\{x_1,x_2,y_1,y_2\}\):
over this set every nonempty \(Y\subseteq\{y_3,\ldots,y_7\}\) with
\(|Y|\leq3\) has \(e_G(Y;\{x_1,x_2,y_1,y_2\})\leq2|Y|\), and
\(|Y|-2|Y|\alpha>0\), so no such \(Y\) is even dense.  (Each of
\(y_3,\ldots,y_6\) has one neighbour in \(\{x_1,x_2,y_1,y_2\}\) and \(y_7\)
has two; a set of at most three of \(y_3,\ldots,y_7\) spans at most three
edges, and at most two if it contains \(y_7\), whose only neighbour among
them is \(y_5\).)  Further
\(\mathrm{cl}_1(x_1,x_2)=\{x_1,x_2\}\), since one vertex adjacent to at most
two vertices of \(\{x_1,x_2\}\) gives \(1-2\alpha>0\), and
\(\mathrm{cl}_0(x_1,x_2)=\{x_1,x_2\}\).  The bound \(t\) thus decides whether
the \(K_5\) is reached.  By \cref{thm:p03-w03-p03-068-well-defined} any
terminal chain gives the closure, so these finite computations suffice.
\end{hbexample}

\section{Closure types as first-order predicates}
\label{sec:p03-w03-p03-068-types}

\begin{hbdefinition}{def:p03-w03-p03-068-closure-iso}{Isomorphic closures; closure types}
Let \(x_1,\ldots,x_r\) be vertices of \(G\) and \(x_1',\ldots,x_r'\)
vertices of \(G'\).  Following Spencer, write
\[
 \mathrm{cl}_t(x_1,\ldots,x_r)\cong\mathrm{cl}_t(x_1',\ldots,x_r')
\]
if there is a graph isomorphism
\[
 G[\mathrm{cl}_t(x_1,\ldots,x_r)]\longrightarrow
 G'[\mathrm{cl}_t(x_1',\ldots,x_r')]
\]
that preserves adjacency and nonadjacency and sends each \(x_i\) to \(x_i'\) \cite[pp.~8--9]{Spencer2007LogicRandomStructuresAuthor}.
For \(t=0\) this says that \(x_i\mapsto x_i'\) is a well-defined graph
isomorphism between the induced graphs on these at most \(r\) vertices
\cite[p.~9]{Spencer2007LogicRandomStructuresAuthor}; in particular \(x_i=x_j\) if and only
if \(x_i'=x_j'\).

A \emph{\(t\)-closure type of arity \(r\)} is a pair \((H,\rho)\) of a finite
graph \(H\) and a map \(\rho\colon\{1,\ldots,r\}\to V(H)\) such that the
\(t\)-closure of \(\rho(\{1,\ldots,r\})\) computed in \(H\) is \(V(H)\).
Write \(\mathrm{cl}_t(x_1,\ldots,x_r)\cong H\) (Spencer's notation, with the
roots understood to be in the positions given by \(\rho\)) if there is an
isomorphism \(\pi\) from \(G[\mathrm{cl}_t(x_1,\ldots,x_r)]\) onto \(H\) with
\(\pi(x_i)=\rho(i)\) for each \(i\) \cite[p.~9]{Spencer2007LogicRandomStructuresAuthor}.
\end{hbdefinition}

By item~6 of \cref{thm:p03-w03-p03-068-well-defined}, whenever
\(\mathrm{cl}_t(x_1,\ldots,x_r)\) is finite the induced graph on it, with the
roots marked, is a \(t\)-closure type; so finite closures are classified by
these types.  Spencer observes that \(\mathrm{cl}_t(x_1,\ldots,x_r)\cong H\) is
a first-order predicate \cite[p.~9]{Spencer2007LogicRandomStructuresAuthor}.

\begin{hbtheorem}{thm:p03-w03-p03-068-type-formula}{proved=annex:p03-w03-p03-068-type-formula}{The closure type is first-order definable}
Let \(t,r\geq0\) and let \((H,\rho)\) be a \(t\)-closure type of arity
\(r\) with \(m=|V(H)\setminus\rho(\{1,\ldots,r\})|\).  There is a
first-order formula \(\varphi_{t,H,\rho}(x_1,\ldots,x_r)\) in the language
\(\{\sim\}\) with equality, of the form
\[
  \exists y_1\cdots\exists y_m\,\bigl[\Delta_{H,\rho}(\bar x,\bar y)\wedge
  \neg\Theta_{t,H}(\bar x,\bar y)\bigr],
\]
such that for every graph \(G\) and all vertices \(x_1,\ldots,x_r\) of \(G\),
\[
  G\entails\varphi_{t,H,\rho}[x_1,\ldots,x_r]
  \iff\mathrm{cl}_t(x_1,\ldots,x_r)\cong H .
\]
Here \(\Delta_{H,\rho}\) is the quantifier-free diagram of \(H\) with the
roots in the positions \(\rho\), and \(\Theta_{t,H}\) is a finite disjunction
of existential formulas, one for each way of attaching at most \(t\) new
vertices rigidly to a copy of \(V(H)\).  For \(t=0\) one has \(m=0\) and
\(\Theta_{0,H}\) is the empty disjunction, so \(\varphi_{0,H,\rho}\) is the
quantifier-free formula \(\Delta_{H,\rho}(\bar x)\).  The formula depends on
\(\alpha\) only through the set of attachments that are rigid.
\end{hbtheorem}

For Spencer's example the formula states the existence of
\(y_1,\ldots,y_7\) with exactly the adjacencies of
\cref{def:p03-w03-p03-068-example}, and, for each of the finitely many rigid
attachments of at most four new vertices to \(x_1,x_2,y_1,\ldots,y_7\), the
nonexistence of vertices \(z_1,\ldots,z_w\) realizing it
\cite[p.~9]{Spencer2007LogicRandomStructuresAuthor}.  The proof also explains why the clause
requiring \(V(H)\) to be the \(t\)-closure of the roots inside \(H\) is
needed: without it the formula would only say that the roots lie in some
\(t\)-closed set of that shape.

\section{The Nonexistence and Finite Closure Lemmas}
\label{sec:p03-w03-p03-068-lemmas}

From here on \(G\sim G(n,p)\) with \(p=n^{-\alpha}\).  For a set \(S\) of
vertices write \(e(S)\) for the number of edges of \(G[S]\), so that
\(e(S)=e_G(S;\emptyset)\).  A nonempty finite \(S\) is \emph{dense} if
\(|S|-e(S)\alpha<0\), that is, if the rooted graph \((\emptyset,G[S])\) is
dense.

\begin{hblemma}{thm:p03-w03-p03-068-no-dense}{proved=annex:p03-w03-p03-068-no-dense}{No dense set of bounded size}
Let \(\alpha>0\) and let \(M\geq1\) be a fixed integer.  Almost surely
\(G(n,n^{-\alpha})\) contains no dense set of at most \(M\) vertices.
\end{hblemma}

\begin{hbtheorem}{thm:p03-w03-p03-068-nonexistence}{proved=annex:p03-w03-p03-068-nonexistence}{Nonexistence Lemma}
For every integer \(t\geq1\), almost surely
\(\mathrm{cl}_t(\emptyset)=\emptyset\) in \(G\sim G(n,n^{-\alpha})\)
\cite[p.~8]{Spencer2007LogicRandomStructuresAuthor}.  Equivalently, almost surely no set of
at most \(t\) vertices is rigid over \(\emptyset\); and for each finite graph
\(H\) with \(v\) vertices and \(e\) edges and \(v-e\alpha<0\), almost surely
\(G\) contains no copy of \(H\).
\end{hbtheorem}

Spencer proves the lemma by the expected number of copies of a dense \(H\),
which is of order \(n^vp^e\) and tends to \(0\), and notes that a fixed \(t\)
leaves only finitely many \(H\) to consider
\cite[p.~8]{Spencer2007LogicRandomStructuresAuthor}.  The annex follows that argument with
the union bound written out.

\begin{hbtheorem}{thm:p03-w03-p03-068-finite-closure}{proved=annex:p03-w03-p03-068-finite-closure}{Finite Closure Lemma}
Let \(\alpha\in(0,1)\) be irrational and \(r,t\geq1\) integers.  Put
\[
  \epsilon=\min\Bigl\{\frac{e\alpha-v}{v}\;:\;v,e\text{ integers},\ 1\leq v\leq t,
  \ e\geq0,\ v-e\alpha\leq0\Bigr\},
  \qquad K=\lceil r/\epsilon\rceil .
\]
Then \(\epsilon\) exists and is positive, and almost surely in
\(G\sim G(n,n^{-\alpha})\)
\[
  |\mathrm{cl}_t(x_1,\ldots,x_r)|<r+K\qquad\text{for all vertices }x_1,\ldots,x_r
\]
\cite[p.~9]{Spencer2007LogicRandomStructuresAuthor}.
\end{hbtheorem}

Spencer's accounting: the \(r\) roots are capital, and each rigid extension
with \(v_i\) new vertices and \(e_i\) new edges costs \(e_i\alpha-v_i\geq
\epsilon v_i\); a closure of size \(r+K\) would therefore be a dense set of
bounded size \cite[p.~9]{Spencer2007LogicRandomStructuresAuthor}.  The annex adjusts one
detail and supplies another.  Spencer describes each step as having fewer
than \(t\) nonroots, while the definition of the closure allows at most
\(t\); and his final strict inequality \(r-K\epsilon<0\), with his choice
\(K=\lceil r/\epsilon\rceil\), needs \(K>r/\epsilon\), which holds because
\(r/\epsilon\) is irrational.  The
tuple \(x_1,\ldots,x_r\) need not consist of distinct vertices.

\begin{hbcorollary}{thm:p03-w03-p03-068-finitely-many-types}{proved=annex:p03-w03-p03-068-finitely-many-types}{Finitely many closure types occur}
Let \(\alpha\in(0,1)\) be irrational and \(r,t\geq1\).  There is a finite set
\(\mathcal C\) of \(t\)-closure types of arity \(r\) such that almost surely
in \(G(n,n^{-\alpha})\), for every \(r\)-tuple \(\bar x\) of vertices,
\(\mathrm{cl}_t(\bar x)\cong H\) for some \((H,\rho)\in\mathcal C\), and so
\(\bar x\) satisfies the disjunction of the formulas \(\varphi_{t,H,\rho}\),
\((H,\rho)\in\mathcal C\), of \cref{thm:p03-w03-p03-068-type-formula}.
\end{hbcorollary}

\begin{hbexample}{def:p03-w03-p03-068-constants}{The constants at \(\alpha=\pi/7\)}
For each \(v\), the smallest \(e\) with \(e\alpha>v\) gives the minimum of
\((e\alpha-v)/v\) at that \(v\).  At \(\alpha=\pi/7\) these are \(e=3,5,7,9\)
for \(v=1,2,3,4\), with values \(0.3464\ldots\), \(0.1219\ldots\),
\((\pi-3)/3=0.0471\ldots\) and \(0.0097\ldots\).  Hence
\(\epsilon=0.3464\ldots\) for \(t=1\), \(0.1219\ldots\) for \(t=2\),
\(0.0471\ldots\) for \(t=3\) and \(0.0097\ldots\) for \(t=4\), and for \(r=2\)
the lemma gives \(K=6,17,43,205\).  So almost surely every pair of vertices
of \(G(n,n^{-\pi/7})\) has a \(4\)-closure of fewer than \(207\) vertices.
The nine-vertex closure of \cref{def:p03-w03-p03-068-example} is well
within this bound; the bound is not claimed to be sharp.
\end{hbexample}

\begin{hbtheorem}{thm:p03-w03-p03-068-t-fixed}{proved=annex:p03-w03-p03-068-t-fixed}{The bound \(t\) must be fixed}
Let \(\alpha\in(0,1)\) be irrational.  Every finite graph containing a dense set contains a
nonempty set rigid over \(\emptyset\).  In \(G\sim G(n,n^{-\alpha})\), almost
surely the whole vertex set is dense, and hence almost surely
\(\mathrm{cl}_n(\emptyset)\neq\emptyset\).
\end{hbtheorem}

\begin{hbwarning}{def:p03-w03-p03-068-lemma-limits}{What the two lemmas do not give}
\begin{itemize}
\item Both lemmas are stated for each fixed \(t\), with \(K\) chosen after
\(r\) and \(t\).  They do not hold with \(t\) growing with \(n\):
by \cref{thm:p03-w03-p03-068-t-fixed}, \(\mathrm{cl}_t(\emptyset)\) is
almost surely nonempty when \(t=n\).  The constant \(K\) of
\cref{thm:p03-w03-p03-068-finite-closure} depends on \(t\) through
\(\epsilon\), as \cref{def:p03-w03-p03-068-constants} shows.
\item The union of the \(t\)-closures over all \(t\) is not a first-order
notion, and in countable models of the almost sure theory a single vertex
can have an infinite union of closures
(\cref{thm:p03-w03-p03-067-sparse-models}).
\item The lemmas are almost sure statements about \(G(n,n^{-\alpha})\).  In
a particular finite graph the closure of \(\emptyset\) can be nonempty: a
copy of \(K_4\) is rigid over \(\emptyset\) when \(\alpha>2/3\), since \(k\) of
its vertices left unnailed carry \(k(7-k)/2\) counted edges and
\(k-\alpha k(7-k)/2<0\) for \(k=1,2,3,4\).
\item Irrationality of \(\alpha\) makes dense and sparse a dichotomy
(\cref{def:p03-w03-p03-065-dense-sparse}).  At rational \(\alpha\) the
quantity \(v-e\alpha\) can vanish, the minimum defining \(\epsilon\) in the
Finite Closure Lemma can be \(0\) (for example \(v=1\), \(e=2\),
\(\alpha=1/2\)), and \(K=\lceil r/\epsilon\rceil\) is then undefined.
\end{itemize}
\end{hbwarning}

\section{A dynamic view: varying the exponent}
\label{sec:p03-w03-p03-068-dynamic}

For fixed irrational \(\alpha\) every first-order sentence holds almost
surely or almost never in \(G(n,n^{-\alpha})\)
(\cref{thm:p03-w03-p03-065-shelah-spencer}).  Spencer now fixes the sentence
and lets \(\alpha\) decrease from \(1\) to \(0\), which roughly follows the
evolution of the random graph as \(p=n^{-\alpha}\) grows
\cite[p.~15]{Spencer2007LogicRandomStructuresAuthor}.

\begin{hbdefinition}{def:p03-w03-p03-068-fa}{The function \(f_A\) and the theories \(T_\alpha\)}
For a first-order sentence \(A\) in the language of graphs and an
irrational \(\alpha\in(0,1)\), let
\[
  f_A(\alpha)=\lim_{n\to\infty}\Pr[G(n,n^{-\alpha})\entails A].
\]
The domain of \(f_A\) is the set of irrational \(\alpha\in(0,1)\); rational
exponents are excluded to avoid their difficulties
\cite[p.~15]{Spencer2007LogicRandomStructuresAuthor}.  Let \(T_\alpha\) be the set of
sentences holding almost surely in \(G(n,n^{-\alpha})\).  By
\cref{thm:p03-w03-p03-065-shelah-spencer} the limit exists and lies in
\(\{0,1\}\), and \(f_A(\alpha)=1\) exactly when \(A\in T_\alpha\)
\cite[p.~15]{Spencer2007LogicRandomStructuresAuthor}; by
\cref{thm:p03-w03-p03-067-as-theory}, \(T_\alpha\) is complete and closed
under logical consequence.
\end{hbdefinition}

Spencer stresses that the theories \(T_\alpha\) have an explicit description,
so that \(f_A\) can be read off without probabilistic calculation
\cite[p.~15]{Spencer2007LogicRandomStructuresAuthor}.  The description, from his Section~5
\cite[p.~10]{Spencer2007LogicRandomStructuresAuthor}, is recorded here in the form the
continuity argument uses.

\begin{hbdefinition}{def:p03-w03-p03-068-axioms}{The axioms of \(T_\alpha\)}
Vertices \(y_1,\ldots,y_v\) are \emph{\(t\)-generic over} \(x_1,\ldots,x_r\)
if, whenever \(z_1,\ldots,z_w\) are distinct from the \(x\)'s and \(y\)'s,
\(w\leq t\), and \(\{z_1,\ldots,z_w\}\) is rigid over
\(\{x_1,\ldots,x_r,y_1,\ldots,y_v\}\), no \(z_i\) is adjacent to any \(y_j\)
\cite[p.~10]{Spencer2007LogicRandomStructuresAuthor}.  The axioms of \(T_\alpha\) are the
sentences of two schemas \cite[p.~10]{Spencer2007LogicRandomStructuresAuthor}.
\begin{itemize}
\item \emph{Nonexistence.}  For each finite graph \(H\) with \(v\) vertices,
\(e\) edges and \(v-e\alpha<0\): there is no copy of \(H\).
\item \emph{Generic Extension.}  For each safe rooted graph \((R,H)\) and each
\(t\geq0\): for all distinct \(x_1,\ldots,x_r\) there are \(y_1,\ldots,y_v\)
that form an \((R,H)\) extension over the \(x\)'s, have no adjacencies
\(y_iy_j\) or \(y_ix_j\) other than those of \(H\), and, if \(t\geq1\), are
\(t\)-generic over the \(x\)'s.
\end{itemize}
Each Generic Extension axiom is a first-order sentence: \(t\)-genericity is
the conjunction, over the finitely many ways of attaching \(w\leq t\) new
vertices to \(r+v\) given ones that are rigid and have an edge to some
\(y_j\), of the nonexistence of such new vertices.  (Spencer writes the
roots in items 1 and 3 of the schema as \(x_1,\ldots,x_v\); the roots are
\(x_1,\ldots,x_r\).  He also writes the outer quantifier without the word
``distinct''; it is read here, as for \(\operatorname{Ext}(R,H)\) in
\cref{def:p03-w03-p03-065-rooted-graph}, over distinct roots, since an
\((R,H)\) extension is defined over distinct roots.)
\end{hbdefinition}

\begin{hbtheorem}{thm:p03-w03-p03-068-axiomatized}{proved=annex:p03-w03-p03-068-axiomatized}{Almost sure theory, axiomatized form}
Let \(\alpha\in(0,1)\) be irrational.  Every axiom of
\cref{def:p03-w03-p03-068-axioms} holds almost surely in
\(G(n,n^{-\alpha})\).  Together with the graph axioms they form a
consistent complete theory; \(T_\alpha\) is exactly the set of its logical
consequences
\cite[pp.~10--12]{Spencer2007LogicRandomStructuresAuthor}.
\end{hbtheorem}

The source proves the Nonexistence schema by the Nonexistence Lemma, sketches
Generic Extension from a large-deviation estimate that it does not prove, and
proves completeness by an Ehrenfeucht game whose Duplicator strategy uses
\(t\)-closures, the Finite Closure Lemma and closure types
\cite[pp.~10--12]{Spencer2007LogicRandomStructuresAuthor}.  The handbook
supplies the probabilistic estimates in
\cref{thm:p03-w03-p03-064-large-deviation,thm:p03-w03-p03-064-theory-as}
and the closure strategy in \cref{thm:p03-w03-p03-070-extension}.
\Cref{thm:p03-w03-p03-067-axioms-generate} identifies the present schemas
with those used in those proofs.  The annex uses that identification to
derive the axiomatized statement here, so the continuity argument below
has a proved dependency within the handbook.

\begin{hbtheorem}{thm:p03-w03-p03-068-continuity}{proved=annex:p03-w03-p03-068-continuity}{Continuity of \(f_A\)}
Let \(A\) be a first-order sentence and \(\alpha_0\in(0,1)\) irrational.
There is \(\delta>0\) such that \(f_A(\alpha)=f_A(\alpha_0)\) for every
irrational \(\alpha\) with \(|\alpha-\alpha_0|<\delta\)
\cite[p.~15]{Spencer2007LogicRandomStructuresAuthor}.
\end{hbtheorem}

Spencer's argument: \(A\) (or \(\neg A\)) follows from finitely many axioms of
\(T_{\alpha_0}\), which depend on \(\alpha_0\) only through the signs of
finitely many numbers \(v-e\alpha_0\); each sign persists on an interval, so
the same axioms belong to \(T_\alpha\) for \(\alpha\) near \(\alpha_0\)
\cite[p.~15]{Spencer2007LogicRandomStructuresAuthor}.

\begin{hbtheorem}{thm:p03-w03-p03-068-distinct}{proved=annex:p03-w03-p03-068-distinct}{The theories \(T_\alpha\) are pairwise distinct}
For every rational \(q\in(0,1)\) there is a rooted graph \((R_q,H_q)\) such
that, for irrational \(\alpha\in(0,1)\), \(\operatorname{Ext}(R_q,H_q)\) holds
almost surely in \(G(n,n^{-\alpha})\) if \(\alpha<q\) and almost never if
\(\alpha>q\).  Consequently \(T_\alpha\neq T_\beta\) for all irrational
\(\alpha<\beta\) in \((0,1)\).
\end{hbtheorem}

The rooted graph is a cycle of nonroots, each joined to a few roots, with
the extra root edges spread evenly around the cycle; it is given in the
annex.  Spencer states the distinctness as known and justifies it through a
graph \(H\) whose copies have threshold function \(n^{-a/b}\) for a rational
\(a/b\) between \(\alpha\) and \(\beta\) \cite[p.~15]{Spencer2007LogicRandomStructuresAuthor};
the proof here uses an extension statement in place of that graph, and does
not establish the existence of such an \(H\) for every rational.

\begin{hbdefinition}{def:p03-w03-p03-068-spectrum}{Spectrum}
The \emph{spectrum} \(\operatorname{Sp}(A)\) of a sentence \(A\) is the set of its points
of discontinuity \cite[p.~15]{Spencer2007LogicRandomStructuresAuthor}.  Since \(f_A\) is
defined only on irrationals, this is read as: \(q\in(0,1)\) lies in
\(\operatorname{Sp}(A)\) if and only if, for every \(\delta>0\), \(f_A\) is not constant on
the irrationals of \((q-\delta,q+\delta)\cap(0,1)\).  Since \(f_A\) takes
only the values \(0\) and \(1\) when \(A\) is first-order, this agrees with
the limit formulation of \cref{def:p03-w03-p03-069-spectrum}: a
\(\{0,1\}\)-valued function has a limit at \(q\) through the irrationals
exactly when it is constant on the irrationals of some neighbourhood of
\(q\).
\end{hbdefinition}

By \cref{thm:p03-w03-p03-068-continuity} no irrational lies in \(\operatorname{Sp}(A)\),
so every spectrum is a set of rationals in \((0,1)\)
\cite[p.~15]{Spencer2007LogicRandomStructuresAuthor}.

\begin{hbcorollary}{thm:p03-w03-p03-068-clique-spectra}{proved=annex:p03-w03-p03-068-clique-spectra}{Spectra from cliques}
Let \(A_k\) say that there is a \(K_k\).  Then \(f_{A_k}(\alpha)=1\) for
\(\alpha<2/(k-1)\) and \(f_{A_k}(\alpha)=0\) for \(\alpha>2/(k-1)\).  Hence
\(\operatorname{Sp}(A_4)=\{2/3\}\), \(\operatorname{Sp}(A_5)=\{1/2\}\) and
\(\operatorname{Sp}(A_4\wedge\neg A_5)=\{1/2,2/3\}\) \cite[pp.~15--16]{Spencer2007LogicRandomStructuresAuthor},
while \(\operatorname{Sp}(A_3)=\emptyset\).
\end{hbcorollary}

For \(A_4\wedge\neg A_5\), as \(G\) evolves the probability starts near
\(0\), jumps to \(1\) at \(p=n^{-2/3}\) when copies of \(K_4\) appear, and
falls back to \(0\) at \(p=n^{-1/2}\) when copies of \(K_5\) appear
\cite[p.~16]{Spencer2007LogicRandomStructuresAuthor}.

\begin{hbtheorem}{thm:p03-w03-p03-068-finite-spectra}{proved=annex:p03-w03-p03-068-finite-spectra}{Every finite set of rationals is a spectrum}
For every finite set \(Q\) of rationals in \((0,1)\) there is a first-order
sentence \(A\) with \(\operatorname{Sp}(A)=Q\).
\end{hbtheorem}

Spencer states this, says it needs some technical work, and adds that he
once conjectured that every spectrum is finite, which turned out to be false
\cite[p.~16]{Spencer2007LogicRandomStructuresAuthor}; a sentence with infinite spectrum is the
subject of his Section~9 and of \cref{thm:p03-w03-p03-069-infinite-spectrum}.

\begin{hbmethod}{def:p03-w03-p03-068-method}{Reading a spectrum}
Given a sentence built by Boolean combination from sentences \(B_1,\ldots,B_k\)
whose functions \(f_{B_i}\) are known step functions:
\begin{enumerate}
\item Record for each \(B_i\) the rationals where \(f_{B_i}\) changes value;
for an existence statement \((\emptyset,H)\) or an extension statement these
are among the numbers \(a/e_a\) for the subextensions of the rooted graph
(\cref{sec:p03-w03-p03-065-relationships}).
\item On each open interval between consecutive recorded rationals every
\(f_{B_i}\) is constant, so \(f_A\) is the same Boolean combination of these
constants, as in the proof of \cref{thm:p03-w03-p03-068-finite-spectra}.
\item \(\operatorname{Sp}(A)\) consists of the recorded rationals at which the constant
values of \(f_A\) on the two adjacent intervals differ.
\end{enumerate}
Side condition: the method needs each \(f_{B_i}\) as an explicit step
function with finitely many steps; it says nothing about rational exponents
themselves, where limits may fail to be \(0\) or \(1\)
(\cref{def:p03-w03-p03-065-rational-boundary}).
\end{hbmethod}

\begin{hbwarning}{def:p03-w03-p03-068-dynamic-limits}{Limits of the dynamic view}
\(f_A\) is a function of a fixed exponent; for each \(\alpha\) the limit is
taken with \(\alpha\) held fixed as \(n\to\infty\).  It is silent at rational
\(\alpha\): at \(\alpha=1/2\) the common-neighbour statement has
expected number of witnesses tending to \(1\) and is not governed by the
dense--sparse dichotomy (\cref{def:p03-w03-p03-065-rational-boundary}).  The
continuity theorem gives for each \(\alpha_0\) some interval of constancy but
no bound on its length, and it uses the axiomatized form of the
zero--one law (\cref{thm:p03-w03-p03-068-axiomatized}).  Spectra need not be
finite.
\end{hbwarning}

\section{Relationships}
\label{sec:p03-w03-p03-068-relationships}

Rooted graphs, the dense--sparse dichotomy, rigid and safe rooted graphs,
and the criterion that an extension statement holds almost surely exactly
when its rooted graph is safe are in
\cref{def:p03-w03-p03-065-rooted-graph,def:p03-w03-p03-065-dense-sparse,thm:p03-w03-p03-065-safe-criterion};
the present chapter uses them throughout.  The sparse zero--one law is
\cref{thm:p03-w03-p03-065-shelah-spencer}.  The almost sure theory of any
zero--one law is complete with infinite models
(\cref{thm:p03-w03-p03-067-as-theory}); for \(p=n^{-\alpha}\) its countable
models are not unique, and the union of all \(t\)-closures distinguishes two
of them (\cref{def:p03-w03-p03-067-sparse-notions,thm:p03-w03-p03-067-sparse-models}).
The completeness proof for \(T_\alpha\) chooses its closure bounds by the
Finite Closure Lemma and compares closure types, which is the use for which
\cref{thm:p03-w03-p03-068-finite-closure,thm:p03-w03-p03-068-type-formula}
are designed \cite[p.~11]{Spencer2007LogicRandomStructuresAuthor}.  For constant \(p\) the
zero--one law of \cref{sec:p03-zero-one-law} needs no closure: every extension
axiom holds almost surely, and the limiting probability of a sentence does not
depend on \(p\).

\section{Sources}
\label{sec:p03-w03-p03-068-sources}

All source claims are taken from Spencer's author's version
\cite{Spencer2007LogicRandomStructuresAuthor}: the two definitions of the \(t\)-closure, the
deferred lemmas and the example (p.~8); the Nonexistence Lemma and its proof
(p.~8); isomorphic closures, the case \(t=0\), and the first-order predicate
(pp.~8--9); the Finite Closure Lemma and its proof (p.~9); the dense rooted
graph that is not rigid (p.~7); the omitted lemma that a nonsafe extension
contains a rigid subextension (p.~12); \(t\)-genericity and
the schemas of \(T_\alpha\) (p.~10); the dynamic view, \(f_A\), continuity,
distinctness of the \(T_\alpha\) and the spectrum (p.~15); the clique
examples, finite spectra and the finiteness conjecture (pp.~15--16).  The
proofs of
\cref{thm:p03-w03-p03-068-rigid-local,thm:p03-w03-p03-068-move-base,thm:p03-w03-p03-068-safe-or-rigid,thm:p03-w03-p03-068-well-defined,thm:p03-w03-p03-068-type-formula,thm:p03-w03-p03-068-t-fixed,thm:p03-w03-p03-068-distinct,thm:p03-w03-p03-068-finite-spectra}
are written for this chapter and do not follow a source text.
